% !PRINTED 291
\weberpage{291}
%<<CONT
in which
\weq{4}{\alpha\delta - \beta\gamma \equiv 1 \quad (\mathrm{mod}\ p),}
and if by $\sqrt{\varDelta}$ $(\mathrm{mod}\ p)$ we understand an integer
whose square is congruent with $\varDelta$ with respect to the modulus
$p$, we have only to set:
\[
\begin{aligned}
a &\equiv \alpha\sqrt{\varDelta}, & b &\equiv \beta\sqrt{\varDelta}\\
c &\equiv \gamma\sqrt{\varDelta}, & \partial &\equiv \delta\sqrt{\varDelta}
\end{aligned}
\quad (\mathrm{mod}\ p)
\]
In the form (3) the signs of $\alpha$, $\beta$, $\gamma$, $\delta$ can
still be changed simultaneously.

We ask for those elements $z$ which remain unchanged under a substitution
of the form (3). These are determined by the congruence
\[
z \equiv \frac{\gamma+\delta z}{\alpha+\beta z} \quad (\mathrm{mod}\ p),
\]
or
\weq{5}{\beta z^2 + (\alpha-\delta) z - \gamma \equiv 0 \quad (\mathrm{mod}\ p),}
a congruence which, if it is not identical, has at most two incongruent
roots. To obtain these we write the congruence (5), first under the
assumption that $\beta$ is not divisible by $p$, thus:
\[
\left(\beta z + \frac{\alpha-\delta}{2}\right)^2 \equiv \left(\frac{\alpha-\delta}{2}\right)^2 + \beta\gamma \quad (\mathrm{mod}\ p),
\]
or with the help of (4):
\weq{6}{\left(\beta z + \frac{\alpha-\delta}{2}\right)^2 \equiv \left(\frac{\alpha+\delta}{2}\right)^2 - 1 \quad (\mathrm{mod}\ p).}
Accordingly three different cases are to be distinguished:
\[
\mathrm{I.} \qquad \left(\frac{\alpha+\delta}{2}\right)^2 - 1 \equiv 0 \quad (\mathrm{mod}\ p);
\]
then the congruence (5) has one root; there is one and
only one element $z$ that remains unchanged.
\[
\mathrm{II.} \qquad \left(\frac{\alpha+\delta}{2}\right)^2 - 1 \ \text{\gesperrt{quadratic residue of $p$}};
\]
the congruence (5) has two different roots; there are two
elements $z$ that remain unchanged.
\[
\mathrm{III.} \qquad \left(\frac{\alpha+\delta}{2}\right)^2 - 1 \ \text{\gesperrt{quadratic non-residue of $p$}};
\]
the congruence (5) has no root at all, and all elements $z$
\gesperrt{are displaced.}
